% ch04 -- Die Hyperbit-Idee nach Matzke

\epigraph{Aus \enquote{es gibt Bits} folgen Dimension, drei Farben,
eine Fermion-Generation, der Schwarzschild-Radius.}{Dok.\,326}

\section{Der Ausgangspunkt: ein Vorzeichen}

Matzke beginnt mit $e_i^2 = +1$. Das ist ein \textbf{Bit} --- etwas, das
quadriert Eins ergibt. Kein Raum, keine Zeit, keine Metrik. Nur die
Feststellung: es gibt unterscheidbare Dinge. \status{SETZUNG}

\begin{laierspur}
Stell dir vor, du weißt nur eines: Es gibt Dinge, die sich unterscheiden.
Mehr nicht. Kein Raum, keine Kraft, keine Materie --- nur Unterscheidbarkeit.
Matzkes Hyperbit ist genau das: das mathematische Objekt, das genau diese
Aussage trägt, und nichts weiter.
\end{laierspur}

\begin{matzbox}[Matzke: Axiom und Brücken]
\status{SETZUNG} Einziges Axiom: $e_i^2 = +1$ (Hyperbit).\\[4pt]
Physikalische Einheiten über drei externe Brücken:
\begin{itemize}
  \item Bekenstein: $4\lP^2$ pro Bit (Entropie--Flächen-Relation)
  \item Compton $=$ Schwarzschild: $\lP$ als Längenskala
  \item Landauer--Einstein: $\mbit = \mP \ln 2 / (2\pi)$
\end{itemize}
\end{matzbox}

\section{Warum Cl(6) --- und warum drei Farben}

Nimmt man mehrere $e_i$ und lässt sie antikommutieren, entsteht aus dem
Nichts eine Struktur mit Dimension, Orientierung und Drehungen. $\mathrm{Cl}(6)$
zerfällt in genau drei Witt-Paare --- drei Päckchen aus je einem Erzeuger
und einem Vernichter. \status{B}

Nach Furey ist das genau die Struktur einer Fermion-Generation: drei Farben,
ein Lepton. Die Drei ist nicht gewählt, sie fällt heraus.
\textbf{Geometrie als Buchhaltung über Unterscheidbarkeit} --- nicht als
Vorannahme.

\begin{laierspur}
Du nimmst drei Münzen und legst sie so hin, dass keine die andere berührt.
Schon hast du eine Dreiecksstruktur --- nicht weil du ein Dreieck gezeichnet
hast, sondern weil die Münzen es erzwingen. Genauso: Die Clifford-Algebra
$\mathrm{Cl}(6)$ erzeugt Dreifachstruktur nicht weil man sie hineinbaut,
sondern weil Antikommutativität plus sechs Generatoren nichts anderes
übrig lässt.
\end{laierspur}

\section{Der Schwarzloch-Teil}

Bekenstein: Ein Schwarzes Loch hat Entropie proportional zur Fläche, ein
Bit pro vier Planck-Flächen. Matzke dreht um: Wenn ein Bit eine Fläche ist,
dann ist eine Ansammlung von Bits ein Volumen. Er setzt Compton-wavelength
gleich Schwarzschild-Radius und erhält den Schwarzschild-Radius
\emph{ohne die Einstein-Gleichungen je zu benutzen}. \status{B}

Für $20\,M_\odot$: $59{,}1$\,km. Exakt richtig.

Die Stabilitätsschwelle liegt bei $6{,}41$ Bit. $\mathrm{Cl}(6)$ hat $6{,}000$
Bit. Das Universum besteht aus Fermionen, die exakt $0{,}41$ Bit unterhalb
der Kollapsschwelle sitzen.

\begin{laierspur}
Stell dir vor, du stapelst Lego-Steine. Ab einer bestimmten Höhe
kollabiert der Turm unter seinem eigenen Gewicht. In Matzkes Framework
kollabiert ein System von Bits ab einer bestimmten Bitanzahl zu einem
Schwarzen Loch. Die kritische Zahl: 6,41. Ein Fermion hat 6,00 Bits ---
gerade knapp darunter.
\end{laierspur}

\section{Wo die Eleganz erkauft ist}

Die $6{,}41$ hängt an einer postulate: Um von Bits zu Massen zu kommen,
braucht Matzke einen Umrechnungskurs. Er nimmt Landauer ($E = k_BT\ln 2$)
und wertet bei Planck-Temperatur aus:
\[
  \mbit = \frac{\mP \ln 2}{2\pi} \approx 0{,}11\,\mP
\]
Eine \emph{universelle} Bit-Masse. \status{SETZUNG}

Nur: Landauer ist eine Aussage über thermische carrier-Ensembles, nicht
über abstrakte Information (A271--A273). Und die Planck-Temperatur ist
nicht die Temperatur irgendeines physikalischen Systems --- sie ist eine
Energieskala. Die Wahl $T = T_P$ ist eine postulate, die sich nicht selbst
begründet.

\begin{keyresult}[Kürzester Vergleich]
Matzkes Algebra ist schön, weil sie zeigt, wie weit man mit reiner
Kombinatorik kommt. Sie ist \emph{einfach}, weil sie an drei Stellen
aufhört zu fragen --- drei geliehene Brücken.\\[4pt]
FFGFT zahlt den Preis vorne ($\xi$ ableiten, $L_0$, den torus
rechtfertigen) und braucht danach keine Brücken. Beide Rechnungen
gehen auf --- sie sehen an verschiedenen Stellen elegant aus.
\end{keyresult}
